Throughout, the frame is the constant frame $\gamma_{0}, \gamma_{i}$ of the flat background and a dot is $\partial_{t}$, $t = x \cdot \gamma_{0}$. Signs are the store's (de Sitter: $\mathcal{R}(B) = \tfrac{\Lambda}{3}B$).

\emph{(i).} The torsion equation for $\hb^{-1}(\gamma_{i}) = a\gamma_{i}$ gives $\Omega(\gamma_{0}) = 0$ and $\Omega(\hb(\gamma_{i})) = -H\,\gamma_{0} \wedge \gamma_{i}$, and the field strength is $\mathcal{R}(\gamma_{0} \wedge \gamma_{i}) = (\dot H + H^{2})\gamma_{0}\gamma_{i} = (\ddot a/a)\gamma_{0}\gamma_{i}$, $\mathcal{R}(\gamma_{i} \wedge \gamma_{j}) = H^{2}\gamma_{i}\gamma_{j}$ (computed on exact jets at four instants). So $\mathcal{R}$ is $A = \ddot a/a$ on the three planes $\gamma_{0}\gamma_{i}$ and $B = H^{2}$ on their duals: the Ricci map is $\mathcal{R}(\gamma_{0}) = 3A\gamma_{0}$, $\mathcal{R}(\gamma_{i}) = (A + 2B)\gamma_{i}$, the scalar $6(A + B)$, the Weyl part zero, and the square of the field strength $12(A^{2} + B^{2})$, so $\mathcal{E} = 12(A^{2} + B^{2}) - 4\bigl(9A^{2} + 3(A + 2B)^{2}\bigr) + 36(A + B)^{2} = 24AB$. For the dual form: $I_{4}$ maps $\gamma_{0}\gamma_{i}$ to the dual plane and squares to $-1$, so $B \mapsto I_{4}\mathcal{R}(I_{4}\mathcal{R}(B))$ is $-AB$ on all six planes and its trace is $-6AB$. On a general field split the field strength as $\mathcal{R} = W + S + Z$: the Weyl part $W$, the scalar part $S = \tfrac{1}{12}\mathcal{R}\,\id$ and the part $Z$ built from the trace-free Ricci map. With $\star$ the duality $B \mapsto I_{4}B$, $\star^{2} = -1$, the parts $W$ and $S$ commute with $\star$ and $Z$ anticommutes with it (the standard duality properties of the three parts; checked exactly on the sample fields). So $(\star\mathcal{R})^{2} = -(W + S - Z)(W + S + Z)$, and an operator that anticommutes with $\star$ has zero trace, which removes the cross terms: $-4\Tr((\star\mathcal{R})^{2}) = 4\Tr((W + S)^{2}) - 4\Tr(Z^{2})$, while the square of the field strength is $4\Tr(\mathcal{R}^{2}) = 4\Tr((W + S)^{2}) + 4\Tr(Z^{2})$. With $4\Tr(W^{2}) = W^{2}$, $4\Tr(S^{2}) = \tfrac16\mathcal{R}^{2}$ and $4\Tr(Z^{2}) = 2\Ric^{2} - \tfrac12\mathcal{R}^{2}$ this is $W^{2} - 2\Ric^{2} + \tfrac23\mathcal{R}^{2} = \mathcal{E}$. It is checked exactly on the Robertson-Walker instants, the radiation field, a vacuum Kasner field and a Bianchi field with both Weyl and Ricci parts. On the radiation field $a = t^{1/2}$: $H = 1/(2t)$, $A = -1/(4t^{2})$, $B = 1/(4t^{2})$, so $\mathcal{R} = 0$, $\mathcal{E} = -\tfrac{24}{16}t^{-4} = -\tfrac32 t^{-4}$ and $\Ric^{2} = 9A^{2} + 3(A + 2B)^{2} = \tfrac34 t^{-4}$, the value of Theorem \claimref{synthesis.radical-normal-tame-wall-is-a-regular-initial-surface-whose-smoothness-selects-no-state}(iv). Where the Weyl part and the scalar both vanish the square of the field strength is $2\Ric^{2}$, so $\mathcal{E} = -2\Ric^{2}$.

\emph{(ii).} By Fact \claimref{synthesis.imported-heat-kernel-coefficients-and-background-field-gauge}(i), without total derivatives the local density of $a_{4}$ is $(4\pi)^{-2}\tfrac{1}{360}\,\mathrm{tr}\bigl(60\mathcal{R}E + 180E^{2} + 5\mathcal{R}^{2} - 2\Ric^{2} + 2\,\mathrm{Riem}^{2} + 30\,\Omega_{ij}\Omega_{ij}\bigr)$, quadratic in the field strength apart from the term $\mathcal{R}E$, in which $E$ is taken with the sign for which the operator is positive on spheres. By hypothesis (A) the anomalous trace of a field is twice the statistics weight $\sigma w$ of clause (iii) of that Fact times this density ($+1$ for a real boson, $-1$ for the squared Dirac operator of a four-component field, $-2$ for the complex ghost). The operators: for a real scalar $E = -\xi\mathcal{R}$ with no connection; for the Dirac field the square of the Dirac operator has $E = -\mathcal{R}/4$ and connection curvature $\tfrac12\mathcal{R}(e_{i} \wedge e_{j})$ acting by left multiplication, and for a gauge vector $E$ is minus the Ricci map on the vector index with connection curvature $v \mapsto \mathcal{R}(e_{i} \wedge e_{j}) \cdot v$, less two minimally coupled ghost scalars (the operators the attachment result uses for the term linear in curvature, \claimref{synthesis.attachment-modules-give-sakharov-weights-cancelling-at-three-separate-generations}(ii); the vector operator is standard and cited from memory). The trace over the four-dimensional spinor module of a product of two bivectors is four times its scalar part, so $\mathrm{tr}\,\Omega_{ij}\Omega_{ij} = \sum_{i \neq j}\br{\mathcal{R}(e_{i} \wedge e_{j})\mathcal{R}(e^{i} \wedge e^{j})}_{0} = -\tfrac12\mathrm{Riem}^{2}$; on the vector index it is $-\mathrm{Riem}^{2}$ (both computed in the algebra on a field with Weyl and Ricci parts). Hence, as coefficients of $(\mathrm{Riem}^{2}, \Ric^{2}, \mathcal{R}^{2})$ in units of $1/(360(4\pi)^{2})$: scalar $(2, -2, 5 - 60\xi + 180\xi^{2})$, which is $(2, -2, 0)$ at $\xi = \tfrac16$; Dirac field $-\bigl(8 - 15,\,-8,\,20 - 60 + 45\bigr) = (7, 8, -5)$; vector $(8 - 30,\,-8 + 180,\,20 - 60) - 2(2, -2, 5) = (-26, 176, -50)$. With $W^{2} = \mathrm{Riem}^{2} - 2\Ric^{2} + \tfrac13\mathcal{R}^{2}$ and $\mathcal{E} = \mathrm{Riem}^{2} - 4\Ric^{2} + \mathcal{R}^{2}$ these are $3W^{2} - \mathcal{E}$, $18W^{2} - 11\mathcal{E}$ (half of it per Weyl field) and $36W^{2} - 62\mathcal{E}$, with no $\mathcal{R}^{2}$ left over. The Weyl-square coefficients $3, 9, 36$ are three times the channel weights $1, 3, 12$ that \claimref{synthesis.gravitational-pair-production-from-the-operator-source-at-first-order}(iv) computes from pair sources; that both are the weight of the field in the trace-free part of the two-point function of its stress is standard and cited from memory, and the match is used only as a test. The stored coefficient fixes the size of the Euler term and the relative signs of the three spins; the overall sign with which it enters the trace $\rho - 3P$ in the store's signature is an input of (A), taken from Mottola and Vaulin, arXiv:gr-qc/0604051: their Eq. (3.6b) gives the coefficient of the Euler density in the trace as $-(N_{S} + 11N_{F} + 62N_{V})/(360(4\pi)^{2})$ in the signature $(-,+,+,+)$, $N_{F}$ counting Dirac fields, and their Eq. (4.12) gives on the de Sitter field a stress with density $+H^{4}(N_{S} + 11N_{F} + 62N_{V})/(960\pi^{2})$; in the store's signature, for one conformally coupled scalar, the density is $H^{4}/(960\pi^{2})$ and the trace is $4\rho = 24H^{4}/(5760\pi^{2}) = \alpha\mathcal{E}$. So the trace is $+\alpha\mathcal{E}$ with $\alpha$ the sum of the Euler weights over $360(4\pi)^{2} = 5760\pi^{2}$. The content (\claimref{synthesis.krein-quartet-drops-from-the-induced-einstein-hilbert-coefficient-only-under-a-regulator-that-treats-it-alike}(iv)): $48 \cdot \tfrac{11}{2} + 12 \cdot 62 = 264 + 744 = 1008$, and $1016$ with eight real scalars; $1016/5760 = 127/720$. On a field with $\mathcal{R} = 0$ the remainder $5\mathcal{R}^{2}$ of a minimally coupled scalar is zero, so the count $1016$ does not depend on the doublets' curvature coupling there; a scalar that is not conformally coupled has in addition a state-dependent trace, outside (A). As a test against the store: with $n_{1/2} = 4n_{1}$ and $n'_{0} = 3n_{1}$ (Fact \claimref{synthesis.imported-sakharov-induced-gravity-coefficients}(c)), both coefficients cancel exactly when each dimension-zero scalar weighs $-(4 \cdot \tfrac{11}{2} + 62)/3 = -28$ and $-(4 \cdot 9 + 36)/(3 \cdot 3) = -8$ scalar units, the published values (from memory).

\emph{(iii).} The expected stress is conserved (Fact \claimref{synthesis.imported-renormalised-stress-tensor-of-hadamard-dirac-states} for the Dirac field; part of (A) for the others). Conservation of a stress of the form $T(\gamma_{0}) = \rho\gamma_{0}$, $T(\gamma_{i}) = -P\gamma_{i}$ on the Robertson-Walker field is $\dot\rho + 3H(\rho + P) = 0$; with $\rho - 3P = \Theta$ it reads $\dot\rho + 4H\rho = H\Theta$, that is $\tfrac{d}{dt}(a^{4}\rho) = a^{4}H\Theta$. For $\Theta = \alpha\mathcal{E} = 24\alpha H^{2}(\dot H + H^{2})$: $\tfrac{d}{dt}(a^{4}H^{4}) = 4a^{4}H^{3}(\dot H + H^{2})$, so $a^{4}H\Theta = 6\alpha\,\tfrac{d}{dt}(a^{4}H^{4})$. For $\Theta = b\,\Box\mathcal{R} = b(\ddot{\mathcal{R}} + 3H\dot{\mathcal{R}})$, put $\rho_{b} = b(H\dot{\mathcal{R}} - \tfrac1{12}\mathcal{R}^{2} + H^{2}\mathcal{R})$ and use $\dot H = \tfrac16\mathcal{R} - 2H^{2}$: $\dot\rho_{b} + 4H\rho_{b} = b\bigl[(\tfrac16\mathcal{R} - 2H^{2})\dot{\mathcal{R}} + H\ddot{\mathcal{R}} - \tfrac16\mathcal{R}\dot{\mathcal{R}} + 2H(\tfrac16\mathcal{R} - 2H^{2})\mathcal{R} + H^{2}\dot{\mathcal{R}} + 4H^{2}\dot{\mathcal{R}} - \tfrac13H\mathcal{R}^{2} + 4H^{3}\mathcal{R}\bigr] = bH(\ddot{\mathcal{R}} + 3H\dot{\mathcal{R}})$. The general solution of the linear equation is these two plus $C/a^{4}$. (Both identities are also checked as identities in the Taylor coefficients of an arbitrary scale factor.) The map $T_{A}$: with the Ricci eigenvalues $r_{0} = 3A$, $r_{1} = A + 2B$ of (i), $\tfrac{1}{2\alpha}T_{A}(\gamma_{0}) = \bigl[r_{0}^{2} - \tfrac23\mathcal{R}r_{0} - \tfrac12(r_{0}^{2} + 3r_{1}^{2}) + \tfrac14\mathcal{R}^{2}\bigr]\gamma_{0} = 3B^{2}\gamma_{0}$, so $T_{A}(\gamma_{0}) = 6\alpha H^{4}\gamma_{0}$; its trace is $2\alpha\bigl(\Ric^{2} - \tfrac23\mathcal{R}^{2} - 2\Ric^{2} + \mathcal{R}^{2}\bigr) = \alpha(-2\Ric^{2} + \tfrac23\mathcal{R}^{2})$, which is $\alpha\mathcal{E}$ wherever the Weyl part vanishes (then the square of the field strength is $2\Ric^{2} - \tfrac13\mathcal{R}^{2}$); the spatial eigenvalue then follows from the trace, $P_{A} = \tfrac13(6\alpha H^{4} - 24\alpha H^{2}\ddot a/a) = \alpha(2H^{4} - 8H^{2}\ddot a/a)$ (computed from the map at four instants). On the radiation field $H^{4} = \tfrac1{16}t^{-4}$ and $H^{2}\ddot a/a = -\tfrac1{16}t^{-4}$: $\rho_{A} = \tfrac38\alpha t^{-4}$, $P_{A} = \tfrac58\alpha t^{-4}$; and $\mathcal{R} \equiv 0$ there, so $\rho_{b}$ and $\Box\mathcal{R}$ vanish.

\emph{(iv).} The trace of $\mathcal{G}(a) = \mathcal{R}(a) - \tfrac12 a\mathcal{R}$ is $-\mathcal{R}$, so (E), the frame equation of Fact \claimref{synthesis.imported-einstein-cartan-matter-coupling}(a) with the expected stress as source, gives $\mathcal{R} = 8\pi G(\rho - 3P)$, which under (A) is $8\pi G(\alpha\mathcal{E} + b\,\Box\mathcal{R})$. On $a = t^{1/2}$ the left side is zero, $\Box\mathcal{R} = 0$, and $\alpha\mathcal{E} = -\tfrac32\alpha t^{-4}$: a contradiction unless $\alpha = 0$, whatever $C$ and $b$. For $a = t^{p}$: $H = p/t$, $\mathcal{R} = 6p(2p - 1)t^{-2}$, $\dot{\mathcal{R}} = -12p(2p - 1)t^{-3}$, so $\rho_{b} = b\bigl[-12p^{2}(2p - 1) - 3p^{2}(2p - 1)^{2} + 6p^{3}(2p - 1)\bigr]t^{-4} = -9b\,p^{2}(2p - 1)t^{-4}$ and the density at order $t^{-4}$ is $\bigl(6\alpha p^{4} - 9b\,p^{2}(2p - 1)\bigr)t^{-4}$, while the left side $3\bar{M}_{P}^{2}H^{2}$ of the equation is of order $t^{-2}$ and $C/a^{4}$ of order $t^{-4p}$, less singular for $p < 1$. So the order $t^{-4}$ must vanish: $2\alpha p^{2} = 3b(2p - 1)$ for $p \neq 0$. (The trace gives the same: $\alpha\mathcal{E} + b\,\Box\mathcal{R} = 12p(p - 1)\bigl[2\alpha p^{2} - 3b(2p - 1)\bigr]t^{-4}$, and conservation relates the two orders by $4(p - 1)\rho_{4} = p\,\Theta_{4}$.) At $p = \tfrac12$ the right side vanishes for every $b$, so the condition is $\alpha = 0$. For $\alpha > 0$ put $q(p) = 2\alpha p^{2} - 6bp + 3b$, a convex quadratic with $q(0) = 3b$, $q(\tfrac12) = \tfrac12\alpha > 0$ and $q(1) = 2\alpha - 3b$. If $b < 0$ then $q(0) < 0 < q(\tfrac12)$: one root is negative and the other lies in $(0, \tfrac12)$. If $b > \tfrac23\alpha$ then $q(\tfrac12) > 0 > q(1)$: one root lies in $(\tfrac12, 1)$ and the other is above $1$, where $C/a^{4}$ is the more singular term and the balance holds only with $C = 0$. If $0 < b < \tfrac23\alpha$ the discriminant $12b(3b - 2\alpha)$ is negative; at $b = 0$ the root is $p = 0$ and at $b = \tfrac23\alpha$ it is the double root $p = 1$, both excluded. So the order $t^{-4}$ balances for a power law with $0 < p < 1$ exactly for $b < 0$ or $b > \tfrac23\alpha$, at $p = \tfrac12 + \alpha/(12b) + O(b^{-2})$ for large $|b|$, and never at $2p = 1$. The balance is necessary for a solution with that leading behaviour; whether such a solution exists is not shown. By Theorem \claimref{synthesis.radical-normal-tame-wall-is-a-regular-initial-surface-whose-smoothness-selects-no-state}(iv) the bracket of $t^{2p}$ is tame in the cosmic coordinate exactly at $2p = 1$.

\emph{(v).} With $b = 0$, (E) on $\gamma_{0}$ is $3\bar{M}_{P}^{2}H^{2} = \rho = C/a^{4} + 6\alpha H^{4}$ (the store's $\mathcal{G}(\gamma_{0}) = -3H^{2}\gamma_{0}$), that is $\tfrac{3\bar{M}_{P}^{4}}{2\alpha}\,x(1 - x) = C/a^{4}$ with $x = 2\alpha H^{2}/\bar{M}_{P}^{2}$. Since $x(1 - x) = \tfrac14 - (x - \tfrac12)^{2} \le \tfrac14$, $C/a^{4} \le 3\bar{M}_{P}^{4}/(8\alpha)$. At late times $C/a^{4} \to 0$ and the radiation law is the root $x \to 0$; by continuity that branch stays in $0 < x < \tfrac12$ until $C/a^{4}$ reaches its largest value, at $a^{4} = 8\alpha C/(3\bar{M}_{P}^{4})$, where $x = \tfrac12$ and $H^{2} = \bar{M}_{P}^{2}/(4\alpha)$. Differentiating, $6\bar{M}_{P}^{2}H(1 - 2x)\dot H = -4HC/a^{4}$, so $\dot H = -2C/\bigl(3\bar{M}_{P}^{2}a^{4}(1 - 2x)\bigr) \to -\infty$ as $x \to \tfrac12$, and $\mathcal{R} = 6(\dot H + 2H^{2})$ is unbounded there. An expanding solution cannot be continued to smaller $a$, where the equation has no real root. With $C = 0$: $x(1 - x) = 0$. (The bounded branch and its end point are in the literature on the anomaly-corrected equation, Fischetti, Hartle and Hu and Anderson, and the root $x = 1$ is Starobinsky's de Sitter solution, $H^{2} = 360\pi/(GN)$ for $N$ scalar units; cited from memory, not used.)

\emph{(vi).} The constant $C$ is identified with the thermal density of the content: $\rho_{r} = C/a^{4} = \tfrac{\pi^{2}}{30}gT^{4}$ with $g = 8 + 2 \cdot 12 + \tfrac78 \cdot 2 \cdot 48 = 116$ for the content of (ii) (the standard density of a free massless gas). $T_{\max}^{4} = \tfrac{30}{\pi^{2}g}\cdot\tfrac{3\bar{M}_{P}^{4}}{8\alpha} = \tfrac{90 \cdot 720}{8 \cdot 127 \cdot 116}\bar{M}_{P}^{4} = \tfrac{2025}{3683}\bar{M}_{P}^{4}$, the factors of $\pi$ cancelling; $0.86^{4} < \tfrac{2025}{3683} < 0.87^{4}$, and with $\bar{M}_{P} = 2.435 \times 10^{18}$ GeV, $T_{\max}$ lies between $2.09$ and $2.12 \times 10^{18}$ GeV. On the radiation field $|\alpha\mathcal{E}| = \tfrac32\alpha t^{-4}$ equals $3\bar{M}_{P}^{2}H^{2} = \tfrac34\bar{M}_{P}^{2}t^{-2}$ at $t^{2} = 2\alpha/\bar{M}_{P}^{2}$; in Planck times $t_{P} = 1/(\sqrt{8\pi}\bar{M}_{P})$, $(t_{\mathrm{eq}}/t_{P})^{2} = 16\pi\alpha = 127/(45\pi)$, between $0.94^{2}$ and $0.95^{2}$. At $T = 6 \times 10^{14}$ GeV, the upper end of the range that \claimref{dark-matter.ceiling-valued-inherited-success} gives for the rung scale under its Hypothesis (U), a mass used here as a temperature: $r = \rho_{r}/\rho_{\max} = (T/T_{\max})^{4}$ lies between $6.6$ and $6.8 \times 10^{-15}$; the small root of $x(1 - x) = r/4$ lies between $r/4$ and $\tfrac{r}{4}(1 + r)$, and the anomalous share of the density is $6\alpha H^{4}/\rho_{r} = x/(1 - x)$, between $1.6$ and $1.7 \times 10^{-15}$.
